% Preámbulo común para las soluciones de los prácticos de Programación 3
\documentclass[11pt,a4paper]{article}

\usepackage[utf8]{inputenc}
\usepackage[T1]{fontenc}
\usepackage{lmodern}
\usepackage[spanish,es-noshorthands]{babel}
\usepackage[margin=2.3cm]{geometry}
\usepackage{amsmath,amssymb,amsthm}
\usepackage{enumitem}
\usepackage{xcolor}
\usepackage{tikz}
\usetikzlibrary{arrows.meta,positioning,shapes,calc}
\usepackage[noend]{algpseudocode}
\usepackage{algorithm}
\usepackage{booktabs}
\usepackage{tcolorbox}
\usepackage[hidelinks]{hyperref}

\floatname{algorithm}{Algoritmo}
\renewcommand{\algorithmicrequire}{\textbf{Entrada:}}
\renewcommand{\algorithmicensure}{\textbf{Salida:}}
\algrenewcommand\algorithmicwhile{\textbf{mientras}}
\algrenewcommand\algorithmicdo{\textbf{hacer}}
\algrenewcommand\algorithmicfor{\textbf{para}}
\algrenewcommand\algorithmicforall{\textbf{para todo}}
\algrenewcommand\algorithmicif{\textbf{si}}
\algrenewcommand\algorithmicthen{\textbf{entonces}}
\algrenewcommand\algorithmicelse{\textbf{sino}}
\algrenewcommand\algorithmicreturn{\textbf{retornar}}
\algrenewcommand\algorithmicfunction{\textbf{función}}
\algrenewcommand\algorithmicprocedure{\textbf{procedimiento}}

\theoremstyle{plain}
\newtheorem{lema}{Lema}[section]
\newtheorem{teorema}[lema]{Teorema}
\newtheorem{proposicion}[lema]{Proposición}
\theoremstyle{definition}
\newtheorem{definicion}[lema]{Definición}
\theoremstyle{remark}
\newtheorem*{observacion}{Observación}

% Caja para el enunciado de cada ejercicio
\newtcolorbox{enunciado}{colback=gray!6,colframe=gray!50,boxrule=0.4pt,arc=2pt,left=6pt,right=6pt,top=4pt,bottom=4pt,fontupper=\small}
% Nota que vincula el ejercicio con parciales/exámenes
\newcommand{\examref}[1]{\par\noindent{\small\color{blue!50!black}\textbf{Relación con parciales/exámenes:} #1}\par\medskip}

\newcommand{\NP}{\mathsf{NP}}
\newcommand{\Pclass}{\mathsf{P}}
\newcommand{\Oh}{\mathcal{O}}

\setlength{\parskip}{0.4em}
\setlength{\parindent}{0pt}

\setlength{\emergencystretch}{2em}

\newcommand{\IS}{\textsc{Independent Set}}
\newcommand{\VC}{\textsc{Vertex Cover}}
\newcommand{\DS}{\textsc{Dominating Set}}
\newcommand{\Cl}{\textsc{Clique}}
\newcommand{\IntS}{\textsc{Interval Scheduling}}
\newcommand{\leP}{\le_P}
\newcommand{\SI}{S\'I}

\title{Práctico 11 --- Reducciones en tiempo polinomial}
\author{Programación 3 -- FING, UdelaR}
\date{}

\begin{document}
\maketitle

\section*{Recordatorio}

Sean $Y$ y $X$ problemas de decisión. Decimos que $Y \leP X$ si toda instancia de $Y$ se puede resolver
con una cantidad polinomial de pasos de cómputo más una cantidad polinomial de llamadas a una caja negra
que resuelve $X$. En todos los ejercicios de este práctico usamos la forma más común (y la que piden las
soluciones oficiales de los parciales): una \emph{transformación} $f$ que, dada una instancia $s$ de $Y$,
construye en tiempo polinomial en $|s|$ una instancia $f(s)$ de $X$ tal que
\[
  s \text{ es instancia \SI{} de } Y \iff f(s) \text{ es instancia \SI{} de } X .
\]
Entonces, para resolver $s$ basta construir $f(s)$ y hacer \emph{una} llamada a la caja negra de $X$.
Para probar la equivalencia demostramos siempre las dos direcciones:
\begin{enumerate}[label=(\roman*)]
  \item si $s$ es \SI{} entonces $f(s)$ es \SI{} (a partir de una solución de $s$ construimos una de $f(s)$);
  \item si $f(s)$ es \SI{} entonces $s$ es \SI{} (a partir de una solución de $f(s)$ construimos una de $s$).
\end{enumerate}
Recordemos además que $\leP$ es transitiva, y que si $Y \leP X$ y $X \in \Pclass$ entonces $Y \in \Pclass$.

%======================================================================
\section{Ejercicio 1}

\begin{enunciado}
\begin{enumerate}[label=(\alph*)]
  \item Consideramos la versión de decisión del problema de planificación de intervalos (\emph{Interval
  Scheduling}, Capítulo 4 de K\&T): dado un conjunto de intervalos en una línea de tiempo, representados
  mediante una lista de pares de enteros, y un entero $k$, ¿existe un subconjunto de al menos $k$
  intervalos que no se solapan entre sí? Determine si \emph{Interval Scheduling} $\leP$ \emph{Vertex Cover}.
  \item Consideramos el problema \emph{Clique}: dado un grafo $G=(V,E)$ representado mediante una matriz
  de adyacencias y un entero $k$, ¿existe $S\subset V$, con al menos $k$ vértices, tal que para todo par de
  vértices $u,v\in S$ existe una arista entre ellos en $G$? Muestre que \emph{Independent Set} $\leP$ \emph{Clique}.
\end{enumerate}
\end{enunciado}

\examref{La parte (a) es exactamente el Segundo parcial 2017, Ej.~4(a) (``¿Es verdad que Programación de
Intervalos $\le_P$ Cubrimiento de Vértices?''), también recopilado en ``Práctica segundo parcial --
soluciones'' (Parcial-2017). La solución oficial responde \emph{Sí}, porque \IntS{} se resuelve en tiempo
polinomial (llamando 0 veces a la caja negra), o alternativamente porque $\IntS\in\NP$ y \VC{} es
NP-completo. La parte (b) del mismo parcial pregunta si $\IS\leP\IntS$: eso es equivalente a $\Pclass=\NP$.}

\subsection*{(a) \IntS{} $\leP$ \VC: \textbf{Sí}}

La respuesta es afirmativa, y no depende de ninguna propiedad particular de \VC: \IntS{} pertenece a
$\Pclass$, y todo problema de $\Pclass$ se reduce polinomialmente a cualquier problema.

\paragraph{\IntS{} está en $\Pclass$.} En el Capítulo~4 se vio que el algoritmo greedy que ordena los
intervalos por tiempo de finalización creciente y elige cada intervalo compatible con el último elegido
obtiene un conjunto de intervalos no solapados de \emph{cardinal máximo}, $\mathrm{OPT}$, en tiempo
$\Oh(n\log n)$. La instancia $(\mathcal I, k)$ es \SI{} si y solo si $\mathrm{OPT}\ge k$.

\paragraph{Reducción.} Damos dos argumentos equivalentes.

\begin{itemize}
  \item \textbf{Como reducción de Cook (caja negra).} Para resolver $(\mathcal I,k)$ ejecutamos el greedy y
  comparamos $\mathrm{OPT}$ con $k$. Esto usa una cantidad polinomial de pasos y \emph{cero} llamadas a la
  caja negra de \VC, lo cual está permitido por la definición de $\leP$.
  \item \textbf{Como transformación de instancias.} Fijamos de antemano dos instancias de \VC:
  \[
    s_{\SI} = (G_1, 0),\ G_1=(\{a\},\emptyset) \qquad\qquad s_{\mathrm{NO}} = (G_2, 0),\ G_2=(\{a,b\},\{(a,b)\}).
  \]
  $s_{\SI}$ es \SI{} (el conjunto vacío cubre las cero aristas de $G_1$) y $s_{\mathrm{NO}}$ es NO (con 0
  vértices no se cubre la arista $(a,b)$). La transformación $f$ ejecuta el greedy sobre $\mathcal I$ y
  devuelve $s_{\SI}$ si $\mathrm{OPT}\ge k$ y $s_{\mathrm{NO}}$ en caso contrario. Es polinomial
  ($\Oh(n\log n)$ más tiempo constante) y, por construcción, $(\mathcal I,k)$ es \SI{} $\iff$ $f(\mathcal I,k)$ es \SI.
\end{itemize}

Otra justificación (también aceptada en la solución oficial): $\IntS\in\Pclass\subseteq\NP$ y \VC{} es
NP-completo, por lo que todo problema de $\NP$, en particular \IntS, se reduce polinomialmente a \VC.

\begin{observacion}
La reducción \emph{inversa}, $\VC\leP\IntS$, no se sabe si vale: si valiera, como $\IntS\in\Pclass$
tendríamos $\VC\in\Pclass$ y, por ser \VC{} NP-completo, $\Pclass=\NP$. Recíprocamente, si $\Pclass=\NP$
la reducción vale con el mismo truco de arriba. Es decir, ``$\VC\leP\IntS$'' es equivalente a
``$\Pclass=\NP$'' (esa es la parte (b) del parcial 2017, con \IS{} en lugar de \VC).
\end{observacion}

\subsection*{(b) \IS{} $\leP$ \Cl}

Recordemos que \IS{} es: dado $G=(V,E)$ y un entero $k$, ¿existe $S\subseteq V$ con $|S|\ge k$ tal que
ningún par de vértices de $S$ es adyacente?

\paragraph{Transformación.} Dada una instancia $(G,k)$ de \IS{} con $G=(V,E)$ y $|V|=n$, construimos la
instancia $(\overline G, k)$ de \Cl, donde $\overline G = (V,\overline E)$ es el \emph{grafo complemento}:
\[
  \overline E = \{(u,v) : u,v\in V,\ u\neq v,\ (u,v)\notin E\}.
\]
El entero $k$ no cambia.

\paragraph{Tiempo polinomial.} Si $M$ es la matriz de adyacencia de $G$ ($n\times n$), la matriz
$\overline M$ de $\overline G$ se obtiene con $\overline M[u][v] = 1 - M[u][v]$ para $u\neq v$ y
$\overline M[u][u]=0$. Esto recorre la matriz una vez: tiempo $\Oh(n^2)$, lineal en el tamaño de la
representación de la instancia. Copiar $k$ es $\Oh(1)$ (o lineal en su cantidad de bits).

\paragraph{Lema clave.} Para todo $S\subseteq V$: $S$ es independiente en $G$ $\iff$ $S$ es un clique en
$\overline G$. En efecto, para $u\ne v$ en $S$, $(u,v)\notin E \iff (u,v)\in\overline E$ por definición de
$\overline E$; la condición ``para todo par'' es la misma en ambos lados.

\paragraph{Equivalencia.}
\begin{enumerate}[label=(\roman*)]
  \item Si $(G,k)$ es \SI{}, existe $S$ independiente en $G$ con $|S|\ge k$. Por el lema, $S$ es un clique
  de $\overline G$ con $|S|\ge k$, así que $(\overline G,k)$ es \SI{} de \Cl.
  \item Si $(\overline G,k)$ es \SI{}, existe un clique $S$ de $\overline G$ con $|S|\ge k$. Por el lema, $S$
  es independiente en $G$, así que $(G,k)$ es \SI{} de \IS.
\end{enumerate}

Por lo tanto, resolvemos \IS{} construyendo $\overline G$ en tiempo $\Oh(n^2)$ y llamando una vez a la caja
negra de \Cl: $\IS\leP\Cl$.

\begin{center}
\begin{tikzpicture}[v/.style={circle,draw,inner sep=1.5pt,minimum size=5mm}, sel/.style={v,fill=blue!20}]
  \begin{scope}
    \node[sel] (a) at (0,1.5) {$a$};
    \node[v] (b) at (1.5,1.5) {$b$};
    \node[sel] (c) at (1.5,0) {$c$};
    \node[v] (d) at (0,0) {$d$};
    \draw (a)--(b)--(c)--(d)--(a);
    \node at (0.75,-0.7) {$G$: $\{a,c\}$ independiente};
  \end{scope}
  \begin{scope}[xshift=5cm]
    \node[sel] (a) at (0,1.5) {$a$};
    \node[v] (b) at (1.5,1.5) {$b$};
    \node[sel] (c) at (1.5,0) {$c$};
    \node[v] (d) at (0,0) {$d$};
    \draw[very thick,blue] (a)--(c);
    \draw (b)--(d);
    \node at (0.75,-0.7) {$\overline G$: $\{a,c\}$ clique};
  \end{scope}
\end{tikzpicture}
\end{center}

\begin{observacion}
Como $\Cl\in\NP$ (certificado: el conjunto $S$; se verifica en $\Oh(n^2)$ que $|S|\ge k$ y que todos los
pares son adyacentes) y \IS{} es NP-completo, esta reducción prueba que \Cl{} es NP-completo. La misma
transformación sirve en el sentido opuesto ($\Cl\leP\IS$), ya que $\overline{\overline G}=G$.
\end{observacion}

%======================================================================
\section{Ejercicio 2}

\begin{enunciado}
Sea $C=\{C_1,\dots,C_k\}$ un conjunto de cláusulas sobre un conjunto de variables booleanas
$X=\{x_1,\dots,x_n\}$, representados mediante listas, donde cada cláusula $C_i$ cumple las siguientes
propiedades:
\begin{itemize}
  \item Es la disyunción de \textbf{a lo sumo} 3 términos (cada término es una variable $x_i$ o su complemento $\overline{x_i}$).
  \item Una misma variable puede aparecer más de una vez en una misma cláusula.
\end{itemize}
Definimos el problema \emph{Al Menos 2-SAT} de la siguiente manera: ¿existen \textbf{al menos 2}
asignaciones de verdad distintas para $X$ que satisfacen $C$? Muestre que \emph{3-SAT} $\leP$ \emph{Al Menos 2-SAT}.
\end{enunciado}

\examref{Misma técnica (agregar una variable nueva ``libre'' o cláusulas auxiliares sobre una variable
nueva) que en el Segundo parcial 2020, Ej.~3 (casiSAT: $\text{SAT}\leP\text{casiSAT}$ agregando $z$ y las
cláusulas $z$ y $\overline z$) y en el Segundo parcial 2024, Ej.~2 ($\text{3-SAT}\leP\text{4-SAT}$
agregando una variable nueva $z$).}

\paragraph{Idea.} Una instancia satisfacible de 3-SAT puede tener una única asignación que la satisface,
así que no basta con copiarla. Agregamos una variable nueva $y$ que no restringe nada: cada asignación
que satisface la instancia original da lugar a \emph{dos} asignaciones distintas de la nueva instancia
(con $y=0$ y con $y=1$).

\paragraph{Transformación.} Sea $(X,C)$ una instancia de 3-SAT, con $X=\{x_1,\dots,x_n\}$ y
$C=\{C_1,\dots,C_k\}$ (cada cláusula con 3 términos). Construimos la instancia $(X',C')$ de Al Menos 2-SAT:
\begin{itemize}
  \item $X' = X\cup\{y\}$, donde $y$ es una variable nueva ($y=x_{n+1}$).
  \item $C' = C\cup\{C_{k+1}\}$ con $C_{k+1} = (y\vee\overline y)$.
\end{itemize}
Las cláusulas de $C$ tienen 3 términos y $C_{k+1}$ tiene 2, de modo que todas cumplen la restricción de
``a lo sumo 3 términos'': $(X',C')$ es una instancia válida. La cláusula $C_{k+1}$ es una tautología (la
satisface cualquier asignación); la agregamos solo para que $y$ aparezca efectivamente en la instancia (en
caso de que las variables se deduzcan de las cláusulas). Si $X$ se da explícitamente como lista, bastaría con
agregar $y$ a esa lista.

\paragraph{Tiempo polinomial.} Copiar las listas de $X$ y $C$ e incorporar una variable y una cláusula de dos
términos lleva tiempo $\Oh(n+k)$, lineal en el tamaño de la instancia de 3-SAT.

\paragraph{Equivalencia.} Para una asignación $\sigma'$ de $X'$ denotamos $\sigma'|_X$ su restricción a $X$.
Observemos que, como las cláusulas $C_1,\dots,C_k$ no mencionan a $y$ y $C_{k+1}$ siempre se satisface,
\begin{equation}\label{eq:restr}
  \sigma' \text{ satisface } C' \iff \sigma'|_X \text{ satisface } C. \tag{$*$}
\end{equation}

\begin{enumerate}[label=(\roman*)]
  \item \emph{Si $(X,C)$ es \SI{} de 3-SAT, entonces $(X',C')$ es \SI{} de Al Menos 2-SAT.} Sea $\sigma$
  una asignación de $X$ que satisface $C$. Definimos $\sigma_0 = \sigma\cup\{y\mapsto 0\}$ y
  $\sigma_1=\sigma\cup\{y\mapsto 1\}$. Ambas son asignaciones de $X'$, son distintas (difieren en $y$) y, por
  \eqref{eq:restr}, ambas satisfacen $C'$ porque su restricción a $X$ es $\sigma$. Hay al menos 2
  asignaciones distintas que satisfacen $C'$.
  \item \emph{Si $(X',C')$ es \SI{} de Al Menos 2-SAT, entonces $(X,C)$ es \SI{} de 3-SAT.} Existen al menos
  dos asignaciones de $X'$ que satisfacen $C'$; en particular existe al menos una, $\sigma'$. Por
  \eqref{eq:restr}, $\sigma'|_X$ satisface $C$, así que $(X,C)$ es satisfacible.
\end{enumerate}

Concluimos que 3-SAT se resuelve construyendo $(X',C')$ en tiempo lineal y haciendo una sola llamada a la
caja negra de Al Menos 2-SAT, o sea $\text{3-SAT}\leP\text{Al Menos 2-SAT}$.

\begin{observacion}
Al Menos 2-SAT está en $\NP$: un certificado son dos asignaciones $\sigma_1,\sigma_2$ (dos tiras de $n$
bits); el certificador comprueba que difieren en algún bit y que cada una satisface las $k$ cláusulas, en
tiempo $\Oh(n+k)$. Junto con la reducción anterior y el hecho de que 3-SAT es NP-completo, se concluye que
Al Menos 2-SAT es NP-completo.
\end{observacion}

%======================================================================
\section{Ejercicio 3}

\begin{enunciado}
(Kleinberg \& Tardos, Ex.~8.2.) Una tienda que trata de analizar el comportamiento de sus clientes a
menudo mantiene una tabla bidimensional $A$, donde las filas corresponden a sus clientes y las columnas
a los productos que vende. La entrada $A[i,j]$ especifica la cantidad del producto $j$ que ha sido comprada
por el cliente $i$. Un ejemplo de tabla se muestra a continuación.
\begin{center}
\begin{tabular}{lcccc}
\toprule
Cliente & Detergente & Cerveza & Pañales & Arena para gatos\\
\midrule
Raj & 0 & 6 & 0 & 3\\
Alanis & 2 & 3 & 0 & 0\\
Chelsea & 0 & 0 & 0 & 7\\
\bottomrule
\end{tabular}
\end{center}
Decimos que un subconjunto $S$ de clientes es \emph{diverso} si ningún par de clientes en $S$ ha comprado
nunca el mismo producto, es decir que para cada producto, a lo sumo un cliente de $S$ lo ha comprado
alguna vez. Un conjunto diverso de clientes resulta útil por ejemplo como grupo objetivo para los estudios
de mercado.

Definimos el problema de \emph{Conjunto Diverso} como sigue: dada una tabla $A$ del tipo definido
anteriormente, de dimensiones $m\times n$, y un número $k\le m$, ¿existe un subconjunto de al menos $k$
clientes que es diverso?

Muestre que \emph{Independent Set} $\leP$ \emph{Conjunto Diverso}.
\end{enunciado}

\examref{No aparece literalmente en los parciales relevados. La reducción es del mismo tipo que la del
Segundo parcial 2018, Ej.~3 (UOPI: $\text{INDSET}\leP\text{UOPI}$, copiando el grafo y fijando pesos 1), y que
la de $\VC\leP\text{CONTRATAR}$ del Segundo parcial 2022, Ej.~3 (vértices $\mapsto$ personas,
aristas $\mapsto$ tareas/productos).}

\paragraph{Idea.} Clientes = vértices, productos = aristas. Cada producto lo compran exactamente los
dos extremos de la arista correspondiente. Dos clientes ``chocan'' (compraron un mismo producto) si y solo si
los vértices son adyacentes.

\paragraph{Transformación.} Sea $(G,k)$ una instancia de \IS{} con $G=(V,E)$, $V=\{v_1,\dots,v_p\}$ y
$E=\{e_1,\dots,e_q\}$.
\begin{itemize}
  \item Si $k>p$, la instancia de \IS{} es trivialmente NO (no hay $k$ vértices). Devolvemos una instancia
  NO fija de Conjunto Diverso que respete $k\le m$, por ejemplo la tabla $2\times 1$ con
  $A=\begin{pmatrix}1\\1\end{pmatrix}$ y $k=2$ (los dos clientes compraron el producto 1, así que no son un
  conjunto diverso). Si $k\le 0$, ambas instancias son trivialmente \SI{} (el conjunto vacío) y la
  transformación de abajo funciona igual (o se devuelve cualquier instancia \SI{} fija).
  \item Si $k\le p$, construimos la tabla $A$ de dimensiones $m\times n$ con $m=p$ (un cliente $c_i$ por
  cada vértice $v_i$) y $n=q$ (un producto $j$ por cada arista $e_j$), definida por
  \[
    A[i,j]=\begin{cases} 1 & \text{si } v_i \text{ es un extremo de } e_j,\\ 0 & \text{en otro caso,}\end{cases}
  \]
  y conservamos el mismo $k$ (que cumple $k\le m$). Si $E=\emptyset$ se puede agregar una columna de ceros para
  no tener una tabla sin columnas; no cambia nada.
\end{itemize}
En otras palabras, $A$ es la matriz de incidencia vértice--arista de $G$.

\paragraph{Tiempo polinomial.} La tabla tiene $p\cdot q$ entradas; se inicializa en 0 en $\Oh(pq)$ y luego,
recorriendo cada arista $e_j=(v_a,v_b)$, se ponen en 1 las entradas $A[a,j]$ y $A[b,j]$. Tiempo total
$\Oh(pq)=\Oh(p^3)$, polinomial en el tamaño de $G$.

\paragraph{Lema clave.} Para $S\subseteq V$, sea $S_c=\{c_i : v_i\in S\}$. Entonces $S$ es independiente en
$G$ $\iff$ $S_c$ es diverso.

\emph{Prueba.} $S_c$ no es diverso $\iff$ existen $c_a\ne c_b$ en $S_c$ y un producto $j$ con
$A[a,j]>0$ y $A[b,j]>0$ $\iff$ (por construcción, la columna $j$ tiene exactamente dos entradas no nulas, en
las filas de los extremos de $e_j$) $e_j=(v_a,v_b)$ $\iff$ existen $v_a\ne v_b$ en $S$ adyacentes $\iff$ $S$
no es independiente. \qed

\paragraph{Equivalencia.} La correspondencia $S\mapsto S_c$ es una biyección entre subconjuntos de vértices y
de clientes que preserva el cardinal.
\begin{enumerate}[label=(\roman*)]
  \item Si $(G,k)$ es \SI{}, existe $S$ independiente con $|S|\ge k$; por el lema, $S_c$ es diverso y
  $|S_c|=|S|\ge k$: la instancia $(A,k)$ es \SI.
  \item Si $(A,k)$ es \SI{}, existe un conjunto diverso $T$ de al menos $k$ clientes; $T=S_c$ para
  $S=\{v_i: c_i\in T\}$, que por el lema es independiente y tiene $|S|=|T|\ge k$: $(G,k)$ es \SI.
\end{enumerate}
(En el caso $k>p$ ambas instancias son NO por construcción.) Por lo tanto $\IS\leP$ Conjunto Diverso.

\paragraph{Ejemplo.} Para el camino $v_1 - v_2 - v_3$ ($e_1=(v_1,v_2)$, $e_2=(v_2,v_3)$) la tabla es
\begin{center}
\begin{tabular}{lcc}
\toprule
Cliente & Producto $e_1$ & Producto $e_2$\\
\midrule
$c_1$ & 1 & 0\\
$c_2$ & 1 & 1\\
$c_3$ & 0 & 1\\
\bottomrule
\end{tabular}
\end{center}
y $\{c_1,c_3\}$ es diverso, como $\{v_1,v_3\}$ es independiente.

\begin{observacion}
Conjunto Diverso está en $\NP$ (certificado: el subconjunto de clientes; se verifica en $\Oh(mn)$ que tiene
al menos $k$ clientes y que cada columna tiene a lo sumo una entrada positiva en esas filas). Con la
reducción anterior, Conjunto Diverso es NP-completo.
\end{observacion}

%======================================================================
\section{Ejercicio 4}

\begin{enunciado}
Sea $G=(V,E)$ un grafo (no dirigido). Decimos que un conjunto de vértices, $S$, es un \emph{conjunto
dominante}, si todo vértice de $V\setminus S$ es adyacente a algún vértice de $S$. Definimos el problema de
decisión \emph{Dominating Set} como: dado un grafo $G$ y un entero $k$, ¿tiene $G$ un conjunto dominante
con a lo sumo $k$ nodos?

\textbf{Observación:} Para $G$ conexo, el hecho de que un conjunto de vértices sea recubridor implica que
es también un conjunto dominante, pero el recíproco no es cierto.

\textbf{Sugerencia:} Dada una instancia $(G,k)$ de \emph{Vertex Cover}, considere la siguiente reducción a
una instancia $(G',k')$ de \emph{Dominating Set}:
\begin{itemize}
  \item $G'$ mantiene todos los nodos y aristas de $G$. Adicionalmente, por cada arista $(u,v)$ de $G$ se
  agrega en $G'$ un nuevo nodo $X_{uv}$ y las dos aristas $(u,X_{uv})$ y $(v,X_{uv})$.
  \item $k'=k+m$, donde $m$ la cantidad de nodos \emph{aislados} de $G$ (nodos de grado 0).
\end{itemize}
\begin{enumerate}[label=(\alph*)]
  \item Muestre que existe una transformación de tiempo polinómico de $G$ en $G'$.
  \item Muestre que si $G$ tiene un VC de a lo sumo $k$ nodos, entonces $G'$ tiene un DS $D$ con
  $|D|\le k+m$ (de a lo sumo $k+m$ nodos).
  \item Muestre que si $G'$ tiene un DS $D$ con $|D|\le k+m$, entonces $G$ tiene un VC de a lo sumo $k$ nodos.

  \textbf{Sugerencia:} Muestre que si $\exists x\in D \mid x\notin G \implies D$ puede transformarse en otro
  DS $D' \mid x'\in G,\ \forall x'\in D' \wedge |D'|\le|D|$.
  \item Muestre que \emph{Vertex Cover} $\leP$ \emph{Dominating Set}.
\end{enumerate}
\end{enunciado}

\examref{No aparece literalmente en los parciales relevados. Es una reducción desde \VC{} con la misma
estructura de prueba que la del Segundo parcial 2022, Ej.~3 ($\VC\leP\text{CONTRATAR}$), cuya redacción
seguimos.}

Notación: $G=(V,E)$, $n=|V|$, $q=|E|$, $I\subseteq V$ el conjunto de vértices aislados de $G$, $m=|I|$.
Llamamos \emph{nodos de arista} a los nodos nuevos $X_{uv}$, y $V'=V\cup\{X_{uv}:(u,v)\in E\}$,
$E'=E\cup\{(u,X_{uv}),(v,X_{uv}) : (u,v)\in E\}$, $G'=(V',E')$.

\begin{center}
\begin{tikzpicture}[v/.style={circle,draw,inner sep=1pt,minimum size=6mm}, xn/.style={rectangle,draw,fill=orange!20,inner sep=2pt,minimum size=5mm}]
  \begin{scope}
    \node[v] (u) at (0,0) {$u$};
    \node[v] (v) at (2,0) {$v$};
    \node[v] (w) at (4,0) {$w$};
    \node[v] (z) at (5.5,0) {$z$};
    \draw (u)--(v)--(w);
    \node at (2.7,-1) {$G$ ($z$ aislado, $m=1$)};
  \end{scope}
  \begin{scope}[xshift=7.5cm]
    \node[v] (u) at (0,0) {$u$};
    \node[v] (v) at (2,0) {$v$};
    \node[v] (w) at (4,0) {$w$};
    \node[v] (z) at (5.5,0) {$z$};
    \node[xn] (xuv) at (1,1.2) {$X_{uv}$};
    \node[xn] (xvw) at (3,1.2) {$X_{vw}$};
    \draw (u)--(v)--(w);
    \draw[orange!70!black] (u)--(xuv)--(v) (v)--(xvw)--(w);
    \node at (2.7,-1) {$G'$};
  \end{scope}
\end{tikzpicture}
\end{center}

Una observación que usaremos varias veces: \emph{todo conjunto dominante de $G'$ contiene a todos los vértices
aislados}, porque un vértice aislado de $G$ sigue aislado en $G'$ (no se agrega ningún $X_{uv}$ adyacente a él)
y por lo tanto, si no estuviera en $D$, no sería adyacente a nadie de $D$. Esto explica el término $+m$ en $k'$:
un VC nunca necesita vértices aislados, pero un DS sí.

\subsection*{(a) La transformación es polinomial}

La construcción de $(G',k')$ consiste en:
\begin{enumerate}
  \item copiar $V$ y $E$: $\Oh(n+q)$;
  \item recorrer $E$ y, por cada arista $(u,v)$, crear el nodo $X_{uv}$ y las aristas $(u,X_{uv})$, $(v,X_{uv})$:
  $\Oh(q)$;
  \item contar los vértices de grado 0 para obtener $m$ (recorriendo las listas de adyacencia, o las filas de la
  matriz): $\Oh(n+q)$ (resp.\ $\Oh(n^2)$), y calcular $k'=k+m$.
\end{enumerate}
$G'$ tiene $n+q$ vértices y $3q$ aristas, de modo que su tamaño es polinomial en el de $G$, y el tiempo total
es $\Oh(n+q)$ con listas de adyacencia, u $\Oh((n+q)^2)$ si se construye la matriz de adyacencia de $G'$. En
cualquier caso, polinomial en el tamaño de la instancia de \VC.

\subsection*{(b) VC de tamaño $\le k$ en $G$ $\Rightarrow$ DS de tamaño $\le k+m$ en $G'$}

Sea $C\subseteq V$ un vertex cover de $G$ con $|C|\le k$. Definimos $D=C\cup I$. Entonces
$|D|\le|C|+|I|\le k+m=k'$. Veamos que $D$ es dominante en $G'$: sea $x\in V'\setminus D$.
\begin{itemize}
  \item Si $x=X_{uv}$ es un nodo de arista: como $C$ cubre la arista $(u,v)$, $u\in C$ o $v\in C$, y ambos son
  adyacentes a $X_{uv}$ en $G'$. Luego $x$ es adyacente a un nodo de $D$.
  \item Si $x\in V$: como $x\notin D\supseteq I$, $x$ no es aislado, así que tiene alguna arista $(x,y)\in E$.
  Como $x\notin C$ y $C$ cubre $(x,y)$, necesariamente $y\in C\subseteq D$, y $(x,y)\in E\subseteq E'$. Luego
  $x$ es adyacente a un nodo de $D$.
\end{itemize}
Así, $D$ es un conjunto dominante de $G'$ con a lo sumo $k+m$ nodos.

\subsection*{(c) DS de tamaño $\le k+m$ en $G'$ $\Rightarrow$ VC de tamaño $\le k$ en $G$}

\begin{lema}[Sugerencia]
Si $D$ es un conjunto dominante de $G'$ que contiene algún nodo de arista, existe un conjunto dominante $D'$
de $G'$ con $D'\subseteq V$ (todos sus nodos están en $G$) y $|D'|\le|D|$.
\end{lema}
\begin{proof}
Tomamos un nodo de arista $X_{uv}\in D$ y definimos $D_1=(D\setminus\{X_{uv}\})\cup\{u\}$. Claramente
$|D_1|\le|D|$ (igual si $u\notin D$, uno menos si ya estaba). Veamos que $D_1$ sigue siendo dominante: los
únicos nodos que podrían haber quedado sin dominar son los que estaban dominados \emph{solo} por $X_{uv}$, o el
propio $X_{uv}$. Los vecinos de $X_{uv}$ en $G'$ son exactamente $u$ y $v$:
\begin{itemize}
  \item $u\in D_1$;
  \item $v$ es adyacente a $u$ en $G'$, pues $(u,v)\in E\subseteq E'$, y $u\in D_1$;
  \item $X_{uv}$ es adyacente a $u\in D_1$.
\end{itemize}
Todo otro nodo no adyacente a $X_{uv}$ sigue dominado por el mismo nodo de $D$ que antes (que sigue en $D_1$).
Luego $D_1$ es dominante y tiene un nodo de arista menos que $D$. Repitiendo el paso para cada nodo de arista
de $D$ (son finitos) se obtiene $D'\subseteq V$ dominante con $|D'|\le|D|$.
\end{proof}

Sea ahora $D$ un conjunto dominante de $G'$ con $|D|\le k+m$. Por el lema podemos suponer (pasando a $D'$)
que $D\subseteq V$, manteniendo $|D|\le k+m$. Por la observación inicial, $I\subseteq D$. Definimos
$C=D\setminus I$. Entonces $|C| = |D|-m\le k$.

Veamos que $C$ es un vertex cover de $G$. Sea $(u,v)\in E$ una arista cualquiera. El nodo $X_{uv}$ no está en
$D$ (porque $D\subseteq V$), así que debe ser adyacente a algún nodo de $D$; sus únicos vecinos son $u$ y $v$,
luego $u\in D$ o $v\in D$. Como $u$ y $v$ tienen grado $\ge1$, no son aislados, así que el que esté en $D$ está
en $C$. Es decir, $C$ cubre $(u,v)$. Como la arista era arbitraria, $C$ es un VC de $G$ con a lo sumo $k$ nodos.

\subsection*{(d) \VC{} $\leP$ \DS}

Dada una instancia $(G,k)$ de \VC, construimos $(G',k')$ como en la sugerencia, en tiempo polinomial por (a).
Por (b), si $(G,k)$ es \SI{} de \VC{} entonces $(G',k')$ es \SI{} de \DS; por (c), si $(G',k')$ es \SI{} de
\DS{} entonces $(G,k)$ es \SI{} de \VC. Por lo tanto
\[
  (G,k) \text{ es \SI{} de \VC} \iff (G',k') \text{ es \SI{} de \DS},
\]
y \VC{} se resuelve con una transformación polinomial más una única llamada a la caja negra de \DS:
$\VC\leP\DS$.

\begin{observacion}
\DS{} está en $\NP$: el certificado es el conjunto $D$ (tira de $|V|$ bits) y el certificador comprueba que
$|D|\le k$ y que cada vértice fuera de $D$ tiene un vecino en $D$, en tiempo $\Oh(|V|+|E|)$. Junto con (d) y
el hecho de que \VC{} es NP-completo, se concluye que \DS{} es NP-completo.

¿Por qué no alcanza con $G'=G$? Porque, como dice la observación del enunciado, un DS no tiene por qué ser
VC: en el triángulo $u,v,w$ el conjunto $\{u\}$ es dominante pero no cubre $(v,w)$. Los nodos $X_{uv}$ obligan
a que cada arista tenga un extremo en el dominante (salvo que se use el propio $X_{uv}$, caso que el lema
elimina).
\end{observacion}

\end{document}
